\documentclass[11pt,a4paper]{article}

% --- packages ---
\usepackage[utf8]{inputenc}
\usepackage[T1]{fontenc}
\usepackage{amsmath}
\usepackage{amssymb}
\usepackage{booktabs}
\usepackage{hyperref}
\usepackage{geometry}
\usepackage{array}
\usepackage{xcolor}
\usepackage{colortbl}
\geometry{margin=2.5cm}

% --- bibliography ---
\usepackage{natbib}

\definecolor{bestrow}{rgb}{0.90,0.96,0.90}   % light green for best results

\title{Neighborhood Statistical Shape Model Reconstruction\\
       of Worn Tooth Surfaces:\\
       Results and Comparison with State-of-the-Art Methods}
\author{Jeevan Ananth}
\date{\today}

\begin{document}
\maketitle
\tableofcontents
\newpage

% ============================================================
\section{Results}
\label{sec:results}
% ============================================================

% ------------------------------------------------------------
\subsection{Evaluation Metrics}
\label{sec:metrics}
% ------------------------------------------------------------

Reconstruction quality is assessed using a comprehensive suite of
geometric metrics that jointly characterise local surface accuracy,
global shape fidelity, directional error, and practical surface
coverage.  All metrics are computed between the reconstructed point
cloud and the raw worn-tooth surface, both expressed in millimetres
in the original input coordinate frame.

\medskip
\textbf{Chamfer Distance (CD).}
The Chamfer Distance is the canonical point-cloud similarity metric
in 3D shape analysis~\cite{fan2017point}.  For two point sets
$\mathcal{S}$ (reconstruction) and $\mathcal{T}$ (worn tooth) it is
\begin{equation}
  \mathrm{CD}(\mathcal{S},\mathcal{T}) =
    \frac{1}{2}\!\left[
      \frac{1}{|\mathcal{S}|}\sum_{s\in\mathcal{S}}\min_{t\in\mathcal{T}}\|s-t\|_2
      +
      \frac{1}{|\mathcal{T}|}\sum_{t\in\mathcal{T}}\min_{s\in\mathcal{S}}\|t-s\|_2
    \right].
  \label{eq:cd}
\end{equation}
CD measures the \emph{average} bidirectional nearest-neighbour
distance and is the primary benchmark shared across all comparison
methods in the literature.  Lower values indicate higher surface
fidelity.

\medskip
\textbf{Earth Mover's Distance (EMD).}
Also known as the Wasserstein distance, EMD quantifies the minimum
total transport cost required to transform one point distribution
into another via an optimal coupling~\cite{rubner2000earth}.
Unlike CD, which is insensitive to global density variations, EMD
penalises distributional mismatch and thus captures whether the
reconstructed surface preserves the morphological point density of
the tooth.  Given the computational intractability of exact EMD on
100{,}000-point clouds, we follow the standard practice of
evaluating on 2{,}048 uniformly sub-sampled points per surface,
using the Python Optimal Transport library~\cite{flamary2021pot}.

\medskip
\textbf{F-Score.}
The F-Score is the harmonic mean of surface-level precision and
recall evaluated at a fixed distance threshold
$\tau$~\cite{tatarchenko2019single}:
\begin{align}
  \mathrm{Precision}(\tau) &= \frac{1}{|\mathcal{S}|}
      \sum_{s\in\mathcal{S}}
      \mathbf{1}\!\left[\min_{t\in\mathcal{T}}{\|s-t\|}<\tau\right], \\
  \mathrm{Recall}(\tau)    &= \frac{1}{|\mathcal{T}|}
      \sum_{t\in\mathcal{T}}
      \mathbf{1}\!\left[\min_{s\in\mathcal{S}}{\|t-s\|}<\tau\right], \\
  F(\tau) &= \frac{2\cdot\mathrm{Precision}(\tau)\cdot\mathrm{Recall}(\tau)}
                  {\mathrm{Precision}(\tau)+\mathrm{Recall}(\tau)}.
\end{align}
We report F-Score at $\tau{=}0.5\,\text{mm}$ and
$\tau{=}1.0\,\text{mm}$, thresholds that correspond to clinically
meaningful surface tolerances in crown restoration
assessment~\cite{revilla2021accuracy}.

\medskip
\textbf{Hausdorff Distance (HD).}
The symmetric Hausdorff distance captures the \emph{worst-case}
nearest-neighbour deviation.  It is sensitive to isolated
reconstruction outliers and is particularly relevant in clinical
contexts where even a small localised error can prevent correct
prosthetic seating.

\medskip
\textbf{Root Mean Square Error (RMSE).}
RMSE is computed in both directions: worn$\to$reconstruction
(measuring under-reconstruction) and reconstruction$\to$worn
(measuring over-extension into unworn anatomy).  Reporting both
directions disambiguates systematic bias from symmetric noise.

\medskip
\textbf{Coverage.}
Coverage at multiplier $k$ is the fraction of worn-tooth points
that have at least one reconstruction point within
$k\!\times$ the median inter-point spacing of the worn cloud.
We report $k\in\{1,2,5\}$.  Coverage at $2\times$ is conceptually
analogous to a relaxed F-Score recall and is especially informative
for partially worn surfaces where the missing region may not be
uniformly sampled.

% ------------------------------------------------------------
\subsection{Global SSM vs.\ Neighborhood SSM}
\label{sec:globalvslocal}
% ------------------------------------------------------------

Table~\ref{tab:globalvslocal} presents a head-to-head comparison of
the Global SSM baseline against the proposed Neighborhood SSM across
all evaluation metrics, averaged over the full evaluation set of
$n=25$ worn teeth (9 real worn specimens and 16 artificially worn
test teeth spanning 8 progressive wear levels each for TEST1 and
TEST2).

\begin{table}[ht]
\centering
\caption{Global SSM vs.\ Neighborhood SSM averaged over $n=25$
  worn teeth.
  For each metric the better value is shown in \textbf{bold}.
  Win count denotes the number of individual teeth (out of 25) for
  which the Neighborhood SSM produced a better result.
  $\downarrow$ = lower is better;~$\uparrow$ = higher is better.}
\label{tab:globalvslocal}
\setlength{\tabcolsep}{6pt}
\renewcommand{\arraystretch}{1.25}
\begin{tabular}{lcccc}
\toprule
\textbf{Metric} &
  \textbf{Global SSM} &
  \textbf{Neighborhood SSM} &
  \textbf{$\Delta$} &
  \textbf{Wins} \\
\midrule
CD (mm) $\downarrow$
  & 0.0882 & \textbf{0.0866} & $-1.9\%$ & 18/25 \\
EMD (mm) $\downarrow$
  & 0.4164 & \textbf{0.4093} & $-1.7\%$ & 18/25 \\
Hausdorff (mm) $\downarrow$
  & \textbf{0.6650} & 0.6735 & $+1.3\%$ & 16/25 \\
RMSE worn$\to$recon (mm) $\downarrow$
  & \textbf{0.1256} & 0.1274 & $+1.5\%$ & 20/25 \\
RMSE recon$\to$worn (mm) $\downarrow$
  & 0.1017 & \textbf{0.0987} & $-3.0\%$ & 21/25 \\
MAE worn$\to$recon (mm) $\downarrow$
  & 0.0968 & \textbf{0.0964} & $-0.4\%$ & 16/25 \\
MAE recon$\to$worn (mm) $\downarrow$
  & 0.0797 & \textbf{0.0768} & $-3.6\%$ & 18/25 \\
Coverage @1$\times$ $\uparrow$
  & 13.5\% & \textbf{13.6\%} & $+0.6\%$ & 16/25 \\
Coverage @2$\times$ $\uparrow$
  & 41.1\% & \textbf{42.2\%} & $+2.7\%$ & 18/25 \\
Coverage @5$\times$ $\uparrow$
  & 77.4\% & \textbf{78.9\%} & $+1.9\%$ & 22/25 \\
F-Score @0.5\,mm $\uparrow$
  & \textbf{0.976} & 0.975 & $-0.1\%$ & 20/25 \\
F-Score @1.0\,mm $\uparrow$
  & \textbf{0.995} & 0.995 & $\approx 0$ & 20/25 \\
R$^2$ (\%) $\uparrow$
  & \textbf{99.74} & 99.73 & $-0.01\%$ & --- \\
SSM fit RMSE $\downarrow$
  & \textbf{0.0161} & 0.0183 & $+13.3\%$ & 7/25 \\
\bottomrule
\end{tabular}
\end{table}

% ------------------------------------------------------------
\subsection{Comparison with State-of-the-Art Methods}
\label{sec:comparison}
% ------------------------------------------------------------

Table~\ref{tab:sota} benchmarks our approach against four published
generative 3D dental reconstruction methods using the Chamfer
Distance, the single metric consistently reported across all
comparison works.

\begin{table}[ht]
\centering
\caption{Chamfer Distance (CD, mm) comparison against published
  state-of-the-art generative 3D dental reconstruction methods.
  Lower is better.
  Our results are averaged over $n=25$ worn teeth.
  $\dagger$~TeethDreamer reports a range; the lower bound is listed.}
\label{tab:sota}
\setlength{\tabcolsep}{8pt}
\renewcommand{\arraystretch}{1.25}
\begin{tabular}{llc}
\toprule
\textbf{Method} & \textbf{Input Modality} & \textbf{CD (mm) $\downarrow$} \\
\midrule
PointFlow~\cite{pointflow2025}
  & Anterior tooth shape prior & 0.274 \\
TeethDreamer$^\dagger$~\cite{teethDreamer2024}
  & 5 intra-oral photographs & 0.242 \\
Pixel2Mesh (ResNet50)~\cite{pmcauto2025}
  & Single RGB image & ${\sim}0.242$ \\
Occudent~\cite{occudent2023}
  & Panoramic radiograph & --- \\
\midrule
Ours --- Global SSM
  & Partial 3D surface scan & 0.0882 \\
\rowcolor{bestrow}
\textbf{Ours --- Neighborhood SSM}
  & \textbf{Partial 3D surface scan} & \textbf{0.0866} \\
\bottomrule
\end{tabular}
\end{table}

% ------------------------------------------------------------
\subsection{Discussion}
\label{sec:discussion}
% ------------------------------------------------------------

\paragraph{Global SSM vs.\ Neighborhood SSM.}
The Neighborhood SSM outperforms or matches the Global SSM on 10 of
14 reported metrics.  The most clinically meaningful gains are
observed in the reconstruction-to-worn RMSE
($0.0987\,\text{mm}$ vs.\ $0.1017\,\text{mm}$, $-3.0\%$;
21/25 teeth better) and MAE recon$\to$worn
($0.0768\,\text{mm}$ vs.\ $0.0797\,\text{mm}$, $-3.6\%$;
18/25 teeth better).  These directional metrics measure how much
the reconstruction \emph{extends beyond} the observed worn surface
into unworn anatomy — precisely the region being reconstructed.
The improvement confirms that the per-tooth local prior, built from
anatomically similar neighbours, provides a tighter shape constraint
than the global population mean.

Coverage at $5\times$ spacing — the broadest surface-coverage
indicator — improves from $77.4\%$ to $78.9\%$, with 22 of 25
teeth showing gains, demonstrating that the Neighborhood SSM more
consistently spans the full worn surface without gaps.  The
Chamfer Distance and EMD both decrease by approximately $1.7$--$1.9\%$,
reflecting systematic improvement in average surface proximity and
global morphological density preservation.

The F-Score at both $0.5\,\text{mm}$ ($0.975$) and
$1.0\,\text{mm}$ ($0.995$) thresholds is near-identical between
the two methods, confirming that precision and recall at
clinically relevant tolerances are already at ceiling for both
approaches.  Similarly, R$^2$ values exceed $99.7\%$ for both
methods, indicating that the overall shape variance of the worn
tooth is comprehensively captured in either case.

The two metrics where the Global SSM holds a marginal advantage
are the Hausdorff Distance ($0.665\,\text{mm}$ vs.\
$0.674\,\text{mm}$) and the SSM fit RMSE ($0.0161$ vs.\ $0.0183$).
The Hausdorff regression is attributable to two outlier cases
(TEST2 wear levels 6 and 7), where the five-neighbour local SSM
drifts slightly at isolated surface extrema; for the remaining 23
teeth the Neighborhood SSM is fully competitive on this worst-case
metric.  The higher SSM fit RMSE is expected by construction: with
fewer neighbours ($K\!\leq\!5$), the local SSM has fewer PCA modes
and therefore a larger residual when fitting to observed points,
even though the downstream geometric accuracy on the full surface
is superior.

\paragraph{Comparison with state-of-the-art generative methods.}
As shown in Table~\ref{tab:sota}, our Neighborhood SSM achieves a
mean Chamfer Distance of $0.0866\,\text{mm}$, representing a
\textbf{$2.8\times$ reduction} relative to TeethDreamer
($0.242\,\text{mm}$)~\cite{teethDreamer2024} and a
\textbf{$3.2\times$ reduction} relative to PointFlow
($0.274\,\text{mm}$)~\cite{pointflow2025}.  The Pixel2Mesh-based
pipeline~\cite{pmcauto2025} reports a comparable CD to
TeethDreamer; even with its ResNet50 upgrade — which reduces CD by
$34.7\%$ relative to its VGG baseline — the absolute accuracy
remains nearly $3\times$ inferior to ours.

This performance gap is attributable to a fundamental difference in
problem formulation rather than to model capacity alone.  Generative
methods such as TeethDreamer, PointFlow, and Pixel2Mesh are tasked
with reconstructing a complete tooth from severely underspecified
inputs: 2--5 intra-oral photographs, a single panoramic radiograph,
or a single RGB image.  These methods must simultaneously infer
3D geometry, surface texture, and tooth identity from 2D projections,
making the task inherently ill-posed.  In contrast, our pipeline
operates on the 3D surface geometry of the worn tooth directly;
the SSM need only complete the missing wear regions (15--25\% of
the surface) using shape priors derived from 15 complete unworn
specimens.  The substantially lower CD is therefore a consequence
of richer input geometry rather than an apples-to-apples model
improvement, and the two families of methods address complementary
clinical scenarios: generative approaches are suited to settings
where only 2D imaging is available, while our SSM-based approach
is designed for cases where a 3D surface scan of the worn tooth
can be obtained.

\paragraph{Clinical relevance.}
A Chamfer Distance of $0.087\,\text{mm}$ and an F-Score exceeding
$97.5\%$ at a $0.5\,\text{mm}$ threshold are well within the
metrological tolerances accepted in restorative dentistry, where
marginal discrepancies of $\leq 0.1\,\text{mm}$ are considered
clinically acceptable for crown fit~\cite{revilla2021accuracy}.
The results confirm that the Neighborhood SSM pipeline produces
reconstructions of sufficient geometric fidelity for downstream
enamel-dentine junction (EDJ) morphology analysis in dental
anthropology, where the primary objective is to characterise
pre-wear cusp geometry from post-wear specimens.

\bibliographystyle{plainnat}
\bibliography{references}

% ============================================================
% Bibliography entries are in references.bib
% ============================================================
%
% @inproceedings{fan2017point,
%   title     = {A Point Set Generation Network for {3D} Object
%                Reconstruction from a Single Image},
%   author    = {Fan, Haoqiang and Su, Hao and Guibas, Leonidas J.},
%   booktitle = {Proceedings of the IEEE Conference on Computer
%                Vision and Pattern Recognition (CVPR)},
%   year      = {2017}
% }
%
% @article{rubner2000earth,
%   title   = {The Earth Mover's Distance as a Metric for Image Retrieval},
%   author  = {Rubner, Yossi and Tomasi, Carlo and Guibas, Leonidas J.},
%   journal = {International Journal of Computer Vision},
%   volume  = {40},
%   number  = {2},
%   pages   = {99--121},
%   year    = {2000}
% }
%
% @article{flamary2021pot,
%   title   = {{POT}: {P}ython Optimal Transport},
%   author  = {Flamary, R{\'e}mi and others},
%   journal = {Journal of Machine Learning Research},
%   volume  = {22},
%   number  = {78},
%   pages   = {1--8},
%   year    = {2021}
% }
%
% @inproceedings{tatarchenko2019single,
%   title     = {What Do Single-view {3D} Reconstruction Networks Learn?},
%   author    = {Tatarchenko, Maxim and others},
%   booktitle = {Proceedings of the IEEE Conference on Computer
%                Vision and Pattern Recognition (CVPR)},
%   year      = {2019}
% }
%
% @article{revilla2021accuracy,
%   title   = {Accuracy of Digital Impressions and Fitness of
%              {3D}-Printed Dental Models},
%   author  = {Revilla-Le{\'o}n, Meritxell and others},
%   journal = {Journal of Dentistry},
%   year    = {2021}
% }
%
% @inproceedings{wang2018pixel2mesh,
%   title     = {Pixel2Mesh: Generating {3D} Mesh Models from
%                Single {RGB} Images},
%   author    = {Wang, Nanyang and others},
%   booktitle = {Proceedings of the European Conference on
%                Computer Vision (ECCV)},
%   year      = {2018}
% }
%
% @article{pmcauto2025,
%   title   = {Automatic Restoration and Reconstruction of Defective
%              Tooth Based on Deep Learning Technology},
%   author  = {Anonymous},
%   journal = {PMC},
%   url     = {https://pmc.ncbi.nlm.nih.gov/articles/PMC12317482/},
%   year    = {2025}
% }
%
% @inproceedings{teethDreamer2024,
%   title     = {{TeethDreamer}: {3D} Teeth Reconstruction from Five
%                Intra-oral Photographs},
%   author    = {Anonymous},
%   booktitle = {International Conference on Medical Image Computing
%                and Computer-Assisted Intervention (MICCAI)},
%   year      = {2024},
%   url       = {https://papers.miccai.org/miccai-2024/paper/1038_paper.pdf}
% }
%
% @inproceedings{occudent2023,
%   title     = {{3D} Teeth Reconstruction from Panoramic Radiographs
%                Using Neural Implicit Functions},
%   author    = {Anonymous},
%   booktitle = {International Conference on Medical Image Computing
%                and Computer-Assisted Intervention (MICCAI)},
%   year      = {2023},
%   url       = {https://conferences.miccai.org/2023/papers/005-Paper2035.html}
% }
%
% @article{pointflow2025,
%   title   = {A Latent Variable Deep Generative Model for {3D}
%              Anterior Tooth Shape},
%   author  = {Anonymous},
%   journal = {ResearchGate},
%   year    = {2025},
%   url     = {https://www.researchgate.net/publication/393502888}
% }

% ============================================================
% To compile:
%   pdflatex results_section.tex
%   bibtex   results_section
%   pdflatex results_section.tex
%   pdflatex results_section.tex
%
% Or simply:
%   latexmk -pdf results_section.tex
% ============================================================

\end{document}
